Over the past two decades, the number of tracked space objects has almost tripled [1]. It is therefore unsurprising to find that low Earth orbit (LEO) has thus become significantly more congested. The operational defence (conjunction screening and collision avoidance manoeuvres) against catastrophic impacts—which risk triggering an irreversible cascade of orbit debris known as the Kessler Syndrome—is now under a significant computational burden [2]. As the catalogue of orbital objects expands so much that traditional mathematical based propagation can no longer keep up—both from an increased number of space objects launched into space each year but also the detection of smaller objects—the need for a new computational method is required [3]. Many organisations that deal with collision avoidance (COLA) such as Advanced Space and GMV, are now turning to machine learning (ML) to solve these problems. This paper explores the efficiency of these ML models in identifying potentially dangerous satellite pairs for further high-fidelity screening just from their orbital characteristics alone, shifting the COLA timeframe further in advance of any potential incidents.
The advent of a new space era, pioneered by mega constellations such as SpaceX’s Starlink, has permanently altered the LEO landscape. Due to private companies’ goals to meet consumer demand, there is sometimes the opportunity cost of the development of COLA architecture. Although there has been no major orbital collision so far in 2026 that has been publicly attributed to failure of conjunction assessment, the operational burden associated with preventing one is rapidly increasing. In a 6-month period (December 2025 – May 2026), Starlink alone conducted more than 200,000 COLA manoeuvres [4].
The central research question is therefore investigating whether a supervised ML model can act as an effective first stage filter for conjunction screening. Rather than replace high-fidelity propagation or the subsequent calculation of collision probability, the proposed model is designed to identify satellite pairs that warrant further analysis at a much earlier stage than is currently industry standard of 7 days for LEO orbits [5]. Pairs classified as significant can then be passed to more computationally expensive screening processes, while those identified as low risk can be deprioritised. This transforms the problem, then. With limited computational resources, the adoption of ML to identify patterns in historical conjunction data significantly reduces the computational workload. This paper studies the trade-offs between accuracy and time in the COLA framework of the NewSpace era, as discussed in ref. [3].
To test this approach, two-line element (TLE) data is taken from historical conjunction data messages (CDMs) and used to construct examples of genuine conjunction events. These are compared against strictly selected negative pairs, representing objects that occupy broadly comparable orbital regions but remain sufficiently separated when propagated. The ML models then use the orbital relationships alone to predict whether the pairs should recall a CDM or not. These models are then evaluated, not only by their overall predictive performance, but also their ability to retain dangerous encounters. This distinction is essential because a missed conjunction is considerably more consequential than an unnecessary secondary evaluation. 
Returning to the central research question, it is hypothesised that the relationships contained within these orbital parameters are sufficiently distinct for an ML model to successfully reduce the number of candidate pairs whilst maintaining a high recall of genuine conjunction events. If supported, this would demonstrate the potential for ML to operate within existing space situational awareness (SSA) architectures, complementing established COLA methods.
